\phantomsection
\subsection*{M3C. Fruit and Spice Mead}
\addcontentsline{toc}{subsection}{M3C. Fruit and Spice Mead}

\textit{Fruit and Spice Mead é um hidromel que contém uma ou mais frutas e uma ou mais especiarias, ervas ou vegetais. Basicamente, trata-se de qualquer hidromel que contenha ao menos um ingrediente especial permitido de cada uma das categorias M2 e M3.}

\textbf{Impressão Geral}: Assemelha-se a um M1 Traditional Mead, porém com um caráter perceptível de fruta e de especiaria, erva ou vegetal (SHV). Uma ampla gama de resultados é possível, dependendo das escolhas do hidromel base e dos ingredientes adicionais, mas o produto final deve ser equilibrado, agradável e prazeroso. O caráter de SHV não deve dominar completamente o hidromel.

\textbf{Aroma}: Os aromas da fruta podem ser percebidos como de um suco doce e fresco até um vinho de fruta envelhecido ou do Porto, dependendo da força alcoólica e do dulçor. Se variedades de fruta forem declaradas, o caráter varietal deve estar perceptível. Se múltiplos tipos de fruta forem declarados, eles não precisam estar em proporção igual, tendo em vista que algumas frutas são mais fortes e distintivas que outras.

Os aromas de SHV podem ser percebidos como de sutis a agressivos, refletindo o caráter dos ingredientes declarados. Muitos descritores são possíveis, como fragrante, aromático, pungente, herbal, picante, floral, frutado, defumado, terroso, resinoso, mentolado, ou termos que remetem a outros ingredientes (semelhante a alcaçuz, a pinho, cítrico, apimentado). Os SHVs declarados devem estar perceptíveis. Se múltiplos SHVs forem declarados, eles não precisam estar em proporção igual, tendo em vista que alguns são mais fortes e distintivos que outros.

O caráter do mel pode variar de sutil a intenso, dependendo da força e dulçor, podendo expressar o néctar floral típico de eventuais variedades de mel declaradas. O buquê de fermentação deve ser equilibrado. Hidroméis mais fortes podem apresentar uma leve nota alcoólica picante, e os carbonatados tendem a expressar mais aroma. Os aromas provenientes de fruta e de SHV devem complementar e realçar o hidromel base, não sobrepujá-lo. Os componentes de fruta, SHV e mel não precisam estar em proporção igual, mas o equilíbrio entre os ingredientes deve ser agradável e prazeroso.

\textbf{Aparência}: De límpido a brilhante; quanto maior a limpidez, mais desejável. A presença de bolhas, espuma, colarinho ou efervescência deve corresponder ao nível de carbonatação declarado. A cor é derivada da variedade de fruta e de mel utilizadas. Quanto maior a vivacidade e o brilho da cor, mais desejável. A cor geralmente não é muito afetada por especiarias, embora flores, pétalas e pimentas possam conferir cores de sutis a significativas; misturas de chá podem conferir cores significativas.

\textbf{Sabor}: Similar ao Aroma em termos de caráter e intensidade de fruta, SHV e mel, tornando-se mais evidente nas versões mais fortes e doces. O caráter de fruta e SHV deve ser perceptível e estar em equilíbrio com o nível de dulçor do hidromel. Algumas frutas e SHVs podem trazer sabores básicos (amargo, doce, azedo, salgado, umami), além de suas características próprias.

Dulçor residual e final de acordo com o nível de dulçor declarado. A acidez deve ser equilibrada, macia ou vibrante, mas não agressiva. Taninos podem fazer um hidromel suave parecer mais seco. Características ásperas, desagradáveis ou excessivamente sulfurosas ou com notas de levedura/autólise são indesejáveis.

\textbf{Sensação na Boca}: A carbonatação é indicada pelo nível declarado, de tranquilo a espumante. O corpo geralmente aumenta com o dulçor e a força alcoólica, tipicamente de médio-leve a cheio. A percepção de álcool acompanha a força alcoólica declarada, variando de ausente a perceptível e com sensação de aquecimento. As adições de fruta e SHV podem aumentar o corpo, aumentar a adstringência ou trazer notas picantes e de aquecimento; nenhum desses aspectos deve ser excessivo. Pimentas podem ser picantes, mas ardência excessiva é um defeito. Qualquer sensação de ardência proveniente da capsaicina deve estar em equilíbrio com o hidromel base e preservar a drinkability.

\textbf{Ingredientes}: Os mesmos ingredientes de um M1 Traditional Mead, com a adição de pelo menos uma fruta e pelo menos um SHV.

\textbf{Comentários}: Se forem utilizadas apenas combinações de ervas ou especiarias, o hidromel deve ser inscrito como M3A, Metheglin. Se forem utilizadas apenas combinações de frutas, deve ser inscrito como M2E Other Fruit Mead. Se forem utilizados outros tipos de ingredientes, deve ser inscrito como M4F Experimental Mead.

Frequentemente, uma combinação de frutas e SHVs pode resultar em um caráter maior do que a soma de suas partes. Os componentes sensoriais provenientes das adições de fruta e SHV devem ser harmoniosos e reforçar os do hidromel base, sem entrar em conflito com eles.

\textbf{Instruções para Inscrição}: Os participantes devem especificar os níveis de dulçor, carbonatação e força alcoólica. Variedades de mel podem ser declaradas. Participantes devem especificar as frutas e os SHVs adicionados. Hidroméis com sabores picantes de pimenta devem especificar o nível de picância (suave, médio, forte) para auxiliar os jurados na ordenação das amostras.

\textbf{Exemplos Comerciais}: Intermiel Pear \& Ginger Sparkling Mead, Lost Cause Banshee Part II, Moonlight Kurt's Apple Pie, Prairie Rose Strawberry Rhubarb, Rabbit's Foot Chocolate Raspberry Love, Schramm's Mead Apple Crisp, St. Ambrose Shotgun Wedding.
